\documentclass[10pt,a4paper]{article}
\usepackage[a4paper,margin=23mm]{geometry}
\usepackage{fontspec}
\setmainfont{Latin Modern Roman}
\setmonofont{lmmono10-regular.otf}[Scale=0.85,BoldFont=lmmonolt10-bold.otf]
\usepackage{amsmath,amssymb,enumitem,listings,xurl}
\usepackage[hidelinks]{hyperref}
\setlength{\parindent}{0pt}
\setlength{\parskip}{4pt}
\setlength{\emergencystretch}{2em}
\setlist{nosep,leftmargin=*}
\lstset{basicstyle=\ttfamily\scriptsize,columns=fullflexible,
breaklines=true,keepspaces=true,showstringspaces=false,aboveskip=6pt,belowskip=6pt}
\title{A Minimal Protocol for Three-Vehicle SITL Validation\\of ROS3-Limassol Receiver-Arrival Independence}
\author{Charalampos Nadiotis\\OpenPELAOS / PELAOS Tech\\
\small Initial project idea: Christos Paparistodemou}
\date{31 October 2026}
\hypersetup{pdftitle={ROS3-Limassol three-vehicle SITL validation protocol},pdfauthor={Charalampos Nadiotis}}
\begin{document}
\maketitle
\vspace{-20pt}
\begin{abstract}
This prospective protocol defines the smallest independently assessable part
of ROS3-Limassol issue~41: two stationary, three-vehicle ArduPilot Rover
software-in-the-loop (SITL) experiments using separate loopback telemetry
paths. Only the second vehicle's path undergoes scripted outages. Existing
receive-only bridges, the one-way outage relay, recording and replay tools
are reused without changing monitor semantics. Acceptance requires detection
and recovery for vehicle~2, absence of stale states for vehicles~1 and~3
during the declared outages, complete evidence records, and live/replay
agreement for all vehicles. This document reports no new experiment or
performance result. Shared-link isolation, authenticated telemetry, physical
radio operation and field validation remain separate evidence requirements.
\end{abstract}

\section{Scope and repository baseline}
Issue~41 distinguishes separate-link setting~A and multiplexed-link
setting~B~\cite{issue41}. This note specifies slice~41A only. Two successful
A runs do not close issue~41, whose acceptance additionally requires two
B runs. A separate path cannot establish the correctness of source filtering
on a common stream.

The inspected repository baseline is commit
\texttt{469ebda}, containing the simulation-paper documentation merge~52;
the relevant runtime and historical evidence baseline is
\texttt{65da1ea17d2a1de68e24a8f6e3120bbfd2fd1e11}.
Issue~33 is closed and PR~38 merged on 27 October 2026. Its twenty
single-vehicle impairment runs are \emph{reported historical SITL results},
not independently reproduced measurements in this work~\cite{pr38,validation}.
The older unchecked live-impairment item in tracking issue~36 must therefore
be interpreted with its later discussion and merged code~\cite{issue36}.

The bridge's default timeout is already 3.0~s. The proposed experiment
explicitly uses 1.0~s at a verified 4~Hz position stream to detect the 2~s
outage; this does not change the default. Issue~41 already has a claimant.
This documentation work does not take over its implementation. Execution
requires coordination with that claimant and a host free of competing SITL
experiments. The earlier different-host exception concerned work before
PR~38 landed; that dependency is now satisfied.

\section{Observation model and proposed method}
For vehicle $v\in\{1,2,3\}$, let $a_v(t)$ be the receiver time of its most
recent accepted telemetry sample at evaluation time $t$. The implemented
monitor rule is
\begin{equation}
L_v(t)=\begin{cases}
\mathrm{WAITING},&\text{no accepted sample exists},\\
\mathrm{STALE},&t-a_v(t)\geq\tau,\\
\mathrm{HEALTHY},&t-a_v(t)<\tau.
\end{cases}
\label{eq:state}
\end{equation}
Here $\tau=1.0$~s and the requested evaluation tick is 0.1~s.
Thus a sample accepted at receiver time 10.0~s becomes stale at the first
evaluation at or after 11.0~s if no later sample is accepted. If the actual
maximum evaluation gap is $\Delta$, detection occurs conditionally within
$[a_v+\tau,a_v+\tau+\Delta]$. A nominal tick is not a measured scheduling
guarantee. Detection relative to outage onset also depends on the final
accepted position sample before the outage.

\newpage
\subsection{Setup and independent perturbation}
Pin a built ArduPilot \texttt{Rover-4.6.3} checkout (historical short commit
\texttt{3fc7011}); record its full SHA and all tool versions. Launch three
stationary instances with instance indices 0--2 and system IDs 1--3, each
with its own MAVProxy process and working directory. Verify the actual
endpoints rather than treating launch arguments as evidence. Expected SITL
TCP ports are 5760/5770/5780. Use UDP~14550 for bridge~1,
14560 for the vehicle~2 relay input, 14611 for bridge~2, and 14570 for
bridge~3. Additional default outputs must not collide or duplicate a
measured endpoint. All measured endpoints bind to \texttt{127.0.0.1}.

Only vehicle~2 traverses \texttt{tools/udp\_outage\_relay.py}.
Its half-open drop windows, relative to its first datagram, are
\begin{equation}
W_1=[30,35)\ \mathrm{s},\qquad W_2=[70,72)\ \mathrm{s}.
\end{equation}
Consequently, the proposed link interruptions last 5~s and 2~s. The relay
remains one-way and single-source; nothing travels from a bridge toward a
vehicle. MAVProxy requests the stream outside the receive-only measurement
path. Before the first window, confirm all three IDs, accepted position
samples, \texttt{SYS\_STATUS}, and stable 4~Hz position delivery. If the
streams have not settled by 30~s, invalidate the attempt rather than
changing a threshold after observing failure.

\subsection{Time alignment and independent controls}
Bridge recordings use independent receiver-relative clocks. For each
vehicle, estimate the wall-clock origin from accepted samples as
\begin{equation}
\widehat{o}_v=\operatorname{median}_k(T_{v,k}-r_{v,k}),
\end{equation}
where $T_{v,k}$ is the sample's receiver wall timestamp expressed in seconds
and $r_{v,k}$ is its recording time in seconds. Convert transitions using
$\widehat{T}_v=\widehat{o}_v+r_v$ and report the spread of the offset
estimates as alignment uncertainty. Apply the relay's wall-clock window
boundaries to all three recordings. Do not assume their start times coincide.

For each control $v\in\{1,3\}$, require
\begin{equation}
\{T:L_v(T)=\mathrm{STALE}\}\cap(W_1\cup W_2)=\varnothing,
\end{equation}
after expressing the windows and states on the same wall-clock axis.
This checks state occupancy, including a stale episode that began before
a window, rather than merely counting transitions occurring inside it.
Retain the wider recovery interval for diagnosis as well.

\section{Acceptance and evidence integrity}
Execute the complete setting twice with fresh filenames (A1 and A2).
For every vehicle~2 run, require two complete relay windows with nonzero
dropped counts, matched stale detection, and subsequent healthy recovery.
Report the measured detection and recovery intervals and alignment spread;
do not substitute a theoretical timing interval for observations.
The battery-age gate may delay accepted recovery samples and must be
explained~\cite{mavlink}.

For all six vehicle recordings, compare ordered live and replay state
sequences after omitting summary records. Require identical sequences and
report maximum absolute transition-time differences. Pass the correct
\texttt{rover1/2/3} robot ID and 1.0~s timeout explicitly. Require zero
\texttt{crc\_errors} and \texttt{bad\_length}; explain all nonzero mapping
and withheld-sample counters. Check startup and deliberate shutdown
separately from declared outage intervals.

\textbf{Immediate stop rule.} If bridge~1 or bridge~3 reacts to vehicle~2's
outage, stop the experiment, retain all evidence and request maintainer
investigation. Treat this as a correctness defect; do not tune monitor
semantics or relax acceptance.

\newpage
PR~51 is open at this protocol's source inspection and addresses incomplete
outage logs being accepted as successful detection~\cite{pr51}. Prefer a
reviewed merge before scoring. Otherwise manually check both complete
start/end/resume record sets, the genuine terminal summary, counters and
recording coverage. A true \texttt{all\_expected\_detected} field alone is
insufficient. Even consistent logs do not prove that all planned time
intervals were observed. Never reconstruct missing evidence to satisfy a
checker.

Retain outside git each original recording, live output, relay log, replay
output and comparison report, together with commands, versions, SHA-256
hashes and an evidence-retention location. Preserve failed attempts.
Do not commit JSONL recordings or logs. After actual execution, add a dated
section to \texttt{robotics/ros3\_limassol/docs/VALIDATION.md}; do not add
unexecuted planned results there.

\section{Implementation boundary and maintainer decisions}
Slice~41A requires no runtime change. Its documentation comprises this note
and a pointer in \texttt{SITL.md}. Run the package unittest suite, demo,
compile checks and \texttt{git diff --check} before the documentation PR;
report unavailable ROS checks as skipped or not run.

The timeout default decision is resolved. MAVLink signing decisions D1--D4
in issue~39 remain unresolved and must precede signing implementation;
signature verification is required before radio-link evaluation~\cite{issue39}.
They do not block unsigned loopback slice~41A. Hardware availability and
paper-publication decisions likewise do not block this protocol.

Future slice~41B adds an optional \texttt{--system-id} selector to the
existing outage relay, keeping whole-datagram dropping unchanged when the
selector is absent. It requires crafted MAVLink~1/2 header tests,
injected-clock boundary tests, single-source checks, and per-window passed
and dropped counts. Verify that every shared output carries all three
vehicles and that all monitored paths experience the same vehicle-2-only
impairment. Every bridge should count other vehicles as
\texttt{other\_source}. The proposed first-header datagram selection is not
general frame-level filtering: if an actual datagram mixes source IDs,
stop for a maintainer scope decision before claiming isolation.

\section{Limitations and reporting status}
This is a prospective protocol, not a new validation result or a submission.
No new SITL experiment was performed for it. Even successful 41A execution
would establish only receiver-arrival behaviour for the tested separate
loopback paths. Shared-link isolation awaits 41B. Physical radio links,
companion-computer deployment, field behaviour, real network impairments,
PX4, signing/authentication, other routers and scalability beyond three
vehicles remain unestablished. \texttt{HEALTHY} describes recent accepted
arrivals, not vehicle safety, source authenticity, capture freshness,
position accuracy or autonomy. This note changes no software or manuscript
licence and leaves the existing manuscript intact.

\appendix
\section{Reproducible proposed commands}
The following Bash commands are an unexecuted run plan. Set absolute paths
to existing checkouts; build pinned Rover SITL using the repository runbook
first~\cite{sitl}. Export the \emph{same resolved} \texttt{AP}, \texttt{PKG}
and \texttt{RUN} values in all launch terminals. Use a fresh directory for A2.
\begin{lstlisting}[language=bash]
export AP=/absolute/path/to/ardupilot
export PKG=/absolute/path/to/ROS3-Limassol/robotics/ros3_limassol
export RUN=$(mktemp -d /tmp/ros3-41A1-XXXXXX)
cd "$PKG"
python -m unittest discover -s tests -v
git diff --check
git rev-parse HEAD
git -C "$AP" rev-parse HEAD
git -C "$AP" describe --tags --always
\end{lstlisting}

\newpage
\textbf{Terminal 1: receivers, started before SITL.}
\begin{lstlisting}[language=bash]
cd "$PKG"
for i in 1 2 3; do
  case "$i" in 1) port=14550;; 2) port=14611;; 3) port=14570;; esac
  python -m ros3_limassol.mavlink_bridge \
    --listen "127.0.0.1:$port" --system-id "$i" --robot-id "rover$i" \
    --stale-after 1.0 --tick 0.1 --duration 180 \
    --record "$RUN/v$i.jsonl" > "$RUN/v$i.out" &
done
wait
\end{lstlisting}
\textbf{Terminal 2: vehicle 2 relay, started before SITL.}
\begin{lstlisting}[language=bash]
cd "$PKG"
python tools/udp_outage_relay.py --listen 127.0.0.1:14560 \
  --forward 127.0.0.1:14611 --outage 30:35 --outage 70:72 \
  --duration 120 --log "$RUN/relay.jsonl" > "$RUN/relay.out"
\end{lstlisting}
\textbf{Terminal 3: three independent stationary instances.}
\begin{lstlisting}[language=bash]
for i in 0 1 2; do
  mkdir -p "$RUN/rover$i"
  (cd "$RUN/rover$i"
   "$AP/Tools/autotest/sim_vehicle.py" -v Rover -I"$i" \
     --sysid "$((i+1))" --no-rebuild \
     --mavproxy-args="--non-interactive --streamrate 4" \
     > "$RUN/sitl$i.out" 2>&1) &
done
wait
\end{lstlisting}
Do not add duplicate default \texttt{--out} endpoints. Record launched
process IDs and terminate only those experiment processes after collection.
The bridge's 180~s timer starts at bridge launch; the relay's 120~s timer
starts at its first datagram. If startup exceeds the 60~s allowance,
invalidate the incomplete attempt and rerun with a longer bridge duration.
Both windows and recovery must lie within each recording.

\textbf{Post-run scoring; inspect evidence completeness first.}
\begin{lstlisting}[language=bash]
cd "$PKG"
python tools/compare_outages.py "$RUN/relay.jsonl" "$RUN/v2.jsonl" \
  --robot-id rover2 --stale-after 1.0 \
  --bridge-output "$RUN/v2.out" --json > "$RUN/v2-compare.json"
for i in 1 2 3; do
  python -m ros3_limassol.replay "$RUN/v$i.jsonl" \
    --robot-id "rover$i" --stale-after 1.0 --json \
    > "$RUN/v$i-replay.jsonl"
done
\end{lstlisting}
The vehicle~2 comparator checks its live/replay transitions. Separately
compare the live and replay sequences for vehicles~1 and~3 and evaluate
their state occupancy on the common wall-clock axis. Do not apply vehicle~2's
outage-detection expectations to the control vehicles.

\begin{thebibliography}{9}
\small
\bibitem{issue41} ROS3-Limassol, issue~41, ``Multi-vehicle SITL: three Rovers,
one bridge per system ID.'' \url{https://github.com/n0711/ROS3-Limassol/issues/41}.
\bibitem{issue36} ROS3-Limassol, tracking issue~36 and maintainer comments.
\url{https://github.com/n0711/ROS3-Limassol/issues/36}.
\bibitem{pr38} ROS3-Limassol, merged PR~38, live impairment on SITL.
\url{https://github.com/n0711/ROS3-Limassol/pull/38}.
\bibitem{validation} ROS3-Limassol, validation record, runtime baseline
\texttt{65da1ea}, \texttt{robotics/ros3\_limassol/docs/VALIDATION.md}.
\bibitem{mavlink} ROS3-Limassol, receive-only MAVLink guide, same baseline,
\texttt{robotics/ros3\_limassol/docs/MAVLINK.md}.
\bibitem{pr51} ROS3-Limassol, PR~51, ``Reject incomplete outage logs before
reporting detection success.'' \url{https://github.com/n0711/ROS3-Limassol/pull/51}.
\bibitem{issue39} ROS3-Limassol, issue~39, signing design for approval.
\url{https://github.com/n0711/ROS3-Limassol/issues/39}.
\bibitem{sitl} ROS3-Limassol, SITL runbook, runtime baseline
\texttt{65da1ea}, \texttt{robotics/ros3\_limassol/docs/SITL.md}.
\end{thebibliography}
\end{document}
